% THEORETICAL FRAMEWORK
% Last generated: 2026-01-15T05:34:34.872Z

We develop a theoretical framework integrating multi-agent systems, mechanism design, and complex systems theory to analyze DAO governance.

\subsection{Formal Model}

\subsubsection{DAO Definition}

A DAO is formally defined as a tuple $\mathcal{D} = (A, T, G, M, S)$ where:

\begin{itemize}
    \item $A = \{a_1, ..., a_n\}$ is the set of agents (members)
    \item $T: A \rightarrow \mathbb{R}^+$ is the token distribution function
    \item $G = (V, Q, \tau)$ is the governance configuration (voting rules, quorum, timelock)
    \item $M$ is the mechanism (how votes are weighted and aggregated)
    \item $S$ is the state (treasury, proposals, delegations)
\end{itemize}

\subsubsection{Agent Model}

Each agent $a_i$ is characterized by:

\begin{equation}
a_i = (\theta_i, \rho_i, \delta_i, \pi_i)
\end{equation}

where:
\begin{itemize}
    \item $\theta_i \in [0, 1]$ is the participation propensity
    \item $\rho_i \in \mathbb{R}^k$ is the preference vector over $k$ proposal dimensions
    \item $\delta_i \in A \cup \{\emptyset\}$ is the delegation target (or self)
    \item $\pi_i$ is the voting strategy function
\end{itemize}

\subsubsection{Proposal Lifecycle}

A proposal $p$ transitions through states:

\begin{equation}
p: \text{Draft} \rightarrow \text{Active} \rightarrow \text{Voting} \rightarrow \{\text{Passed}, \text{Failed}\} \rightarrow \text{Executed}
\end{equation}

The outcome is determined by the aggregation function:

\begin{equation}
\text{outcome}(p) = M\left(\sum_{a_i \in \text{voters}(p)} w(a_i) \cdot v_i(p)\right)
\end{equation}

where $w(a_i)$ is the voting weight and $v_i(p) \in \{0, 1\}$ is the vote (abstention is modeled as non-participation rather than a zero vote).

\subsection{Voting Mechanisms}

\subsubsection{Token-Weighted Voting}

\begin{equation}
w_{\text{token}}(a_i) = T(a_i) + \sum_{a_j: \delta_j = a_i} T(a_j)
\end{equation}

\subsubsection{Quadratic Voting}

\begin{equation}
w_{\text{quad}}(a_i) = \sqrt{T(a_i) + \sum_{a_j: \delta_j = a_i} T(a_j)}
\end{equation}

\subsubsection{Conviction Voting}

Conviction accumulates over time:

\begin{equation}
C_i(t+1) = \alpha \cdot C_i(t) + \left(T(a_i) + \sum_{a_j: \delta_j = a_i} T(a_j)\right) \cdot v_i
\end{equation}

where $\alpha \in (0, 1)$ is the decay factor.

\subsection{Equilibrium Concepts}

\subsubsection{Participation Equilibrium}

An agent participates if expected benefit exceeds cost:

\begin{equation}
\mathbb{E}[U_i(\text{vote})] - c_i > \mathbb{E}[U_i(\text{abstain})]
\end{equation}

This leads to participation equilibria that depend on quorum and voting mechanism.

\subsubsection{Delegation Equilibrium}

Delegation occurs when:

\begin{equation}
U_i(\text{delegate to } a_j) > U_i(\text{vote directly}) - c_{\text{attention}}
\end{equation}

\subsection{Key Theoretical Predictions}

Based on this framework, we derive several testable predictions:

\begin{enumerate}
    \item \textbf{Quorum-Participation Trade-off}: Higher quorums reduce false positives but increase governance failures
    \item \textbf{Scale-Participation Decay}: Participation decreases with DAO size following $\theta(n) \propto n^{-\beta}$
    \item \textbf{Quadratic Egalitarianism}: Quadratic voting reduces Gini coefficient of voting power
    \item \textbf{Delegation Concentration}: Without intervention, delegation concentrates in power-law distributions
    \item \textbf{Pipeline Throughput-Quality Trade-off}: Faster proposal pipelines increase throughput but may reduce decision quality as measured by proposal abandonment rate
    \item \textbf{Treasury Buffer Stabilization}: Emergency reserves and buffer policies reduce treasury volatility at the cost of reduced capital efficiency
    \item \textbf{Cooperation Surplus}: Inter-DAO cooperation generates positive-sum outcomes when participating DAOs have complementary resource profiles
\end{enumerate}

These predictions are tested in Section~\ref{sec:results}.
